% ============================================================================
% STUNDENLEISTUNG - Physik Klasse 6: Bewegungen
% ============================================================================
% Abschlusstest der Unterrichtsreihe "Bewegungen"
% Differenzierung: (*) Basis 26P, (**) Standard 31P, (***) Erweitert 35P
% Dauer: ca. 30 Minuten
% ============================================================================

\documentclass[11pt,a4paper]{article}

\usepackage[utf8]{inputenc}
\usepackage[T1]{fontenc}
\usepackage[ngerman]{babel}
\usepackage{geometry}
\usepackage{helvet}
\usepackage{tikz}
\usepackage{amsmath}
\usepackage{amssymb}
\usepackage{enumitem}
\usepackage{tabularx}
\usepackage{array}
\usepackage{fancyhdr}
\usepackage{lastpage}
\usepackage{xcolor}

\renewcommand{\familydefault}{\sfdefault}

\geometry{left=2cm, right=2cm, top=2.8cm, bottom=2cm}

% Farben
\definecolor{physik}{HTML}{2c5aa0}
\definecolor{lightblue}{HTML}{e3f2fd}
\definecolor{lightgray}{HTML}{f5f5f5}

% Kopfzeile
\pagestyle{fancy}
\fancyhf{}
\fancyhead[L]{\small Physik Klasse 6 | Stundenleistung: Bewegungen}
\fancyhead[R]{\small Name: \underline{\hspace{4cm}} \quad Datum: \underline{\hspace{2.5cm}}}
\fancyfoot[C]{\small Seite \thepage{} von \pageref{LastPage}}
\renewcommand{\headrulewidth}{0.4pt}

% Lücke
\newcommand{\luecke}[1]{\underline{\hspace{#1}}}

% Niveaustufen
\newcommand{\stern}[1]{\textbf{(#1)}}

% Punktefeld
\newcommand{\punkte}[1]{\hfill \textit{/#1 P}}

\begin{document}

% ============================================================================
% TITEL UND HINWEISE
% ============================================================================
\begin{center}
    {\Large\bfseries\textcolor{physik}{Stundenleistung: Bewegungen}}\\[0.3em]
    {\normalsize Physik Klasse 6 | Stunde 4 von 4 | ca. 30 Minuten}
\end{center}

\vspace{0.3em}

\noindent\fcolorbox{physik}{lightblue}{%
\begin{minipage}{\dimexpr\textwidth-2\fboxsep-2\fboxrule}
\textbf{Hinweise:}
\begin{itemize}[leftmargin=1.5em, topsep=2pt, itemsep=0pt]
    \item Bearbeite die Aufgaben in der Reihenfolge \stern{*} $\to$ \stern{**} $\to$ \stern{***}.
    \item Für Diagramme benutze ein Lineal.
    \item Schreibe lesbar und in ganzen Sätzen, wo nötig.
    \item Erlaubte Hilfsmittel: Lineal, Taschenrechner (falls vorhanden).
\end{itemize}
\end{minipage}%
}

\vspace{0.8em}
\hrule
\vspace{0.8em}

% ============================================================================
% AUFGABE 1: Bewegungsarten (*) - 8 Punkte
% ============================================================================
\noindent\textbf{Aufgabe 1 \stern{*}: Bewegungsarten zuordnen} \punkte{8}

\vspace{0.5em}
\noindent Ordne jede Bewegung der richtigen Bahnform und Bewegungsart zu. Setze Kreuze in die Tabelle.

\vspace{0.5em}
\begin{center}
\renewcommand{\arraystretch}{1.4}
\begin{tabular}{|l|c|c|c||c|c|}
\hline
\textbf{Bewegung} & \multicolumn{3}{c||}{\textbf{Bahnform}} & \multicolumn{2}{c|}{\textbf{Bewegungsart}} \\
\cline{2-6}
& geradlinig & kreisförmig & krummlinig & gleichförmig & ungleichförmig \\
\hline
Rolltreppe (konstante Fahrt) & & & & & \\
\hline
Karussell & & & & & \\
\hline
Achterbahn & & & & & \\
\hline
100-m-Sprinter (Start) & & & & & \\
\hline
\end{tabular}
\end{center}

\vspace{1em}

% ============================================================================
% AUFGABE 2: Geschwindigkeit berechnen (*) - 6 Punkte
% ============================================================================
\noindent\textbf{Aufgabe 2 \stern{*}: Geschwindigkeit berechnen} \punkte{6}

\vspace{0.5em}
\noindent Ein Radfahrer legt in 10 Sekunden eine Strecke von 50 Metern zurück.

\vspace{0.3em}
\noindent\textbf{a)} Ergänze die gegebenen und gesuchten Größen:

\vspace{0.3em}
\hspace{1.5em} Gegeben: $s =$ \luecke{2cm} \hspace{2em} $t =$ \luecke{2cm}

\vspace{0.3em}
\hspace{1.5em} Gesucht: \luecke{2cm}

\vspace{0.5em}
\noindent\textbf{b)} Berechne die Geschwindigkeit:

\vspace{0.3em}
\hspace{1.5em} $v = \dfrac{s}{t} = \dfrac{\luecke{2cm}}{\luecke{2cm}} =$ \luecke{2cm}

\vspace{1em}

% ============================================================================
% AUFGABE 3: Werte ablesen (*) - 6 Punkte
% ============================================================================
\noindent\textbf{Aufgabe 3 \stern{*}: Werte aus dem Diagramm ablesen} \punkte{6}

\vspace{0.5em}
\noindent Das Diagramm zeigt die Bewegung eines Fußgängers:

\begin{center}
\begin{tikzpicture}[scale=0.75]
    % Koordinatensystem
    \draw[->] (0,0) -- (7.5,0) node[right] {$t$ in s};
    \draw[->] (0,0) -- (0,5) node[above] {$s$ in m};

    % Skalierung x
    \foreach \x/\label in {0/0, 1/2, 2/4, 3/6, 4/8, 5/10, 6/12} {
        \draw (\x,-0.1) -- (\x,0.1);
        \node[below] at (\x,-0.2) {\small \label};
    }

    % Skalierung y
    \foreach \y/\label in {0/0, 1/10, 2/20, 3/30, 4/40} {
        \draw (-0.1,\y) -- (0.1,\y);
        \node[left] at (-0.2,\y) {\small \label};
    }

    % Gerade (Fußgänger: v = 10/3 m/s ≈ 3,3 m/s → einfacher: 40m in 12s)
    \draw[thick, physik] (0,0) -- (6,4);

    \node[physik, right] at (5.5,3.5) {\small Fußgänger};
\end{tikzpicture}
\end{center}

\noindent Lies die Werte ab:

\begin{enumerate}[label=\alph*), leftmargin=2em]
    \item Welchen Weg hat der Fußgänger nach $t = 6$ s zurückgelegt? $s =$ \luecke{2cm}
    \item Wie lange braucht der Fußgänger für $s = 40$ m? $t =$ \luecke{2cm}
    \item Welchen Weg hat der Fußgänger nach $t = 3$ s zurückgelegt? $s =$ \luecke{2cm}
\end{enumerate}

\vspace{1em}

% ============================================================================
% AUFGABE 4: Diagramm zeichnen (**) - 8 Punkte
% ============================================================================
\noindent\textbf{Aufgabe 4 \stern{**}: Ein s(t)-Diagramm zeichnen} \punkte{8}

\vspace{0.5em}
\noindent Ein Jogger läuft mit gleichförmiger Geschwindigkeit. Seine Messwerte:

\begin{center}
\begin{tabular}{|c|c|c|c|c|c|}
\hline
Zeit $t$ in s & 0 & 4 & 8 & 12 & 16 \\
\hline
Weg $s$ in m & 0 & 20 & 40 & 60 & 80 \\
\hline
\end{tabular}
\end{center}

\noindent\textbf{a)} Zeichne die Punkte in das Koordinatensystem ein und verbinde sie zu einer Geraden.

\begin{center}
\begin{tikzpicture}[scale=0.65]
    % Koordinatensystem
    \draw[->] (0,0) -- (9,0) node[right] {$t$ in s};
    \draw[->] (0,0) -- (0,5.5) node[above] {$s$ in m};

    % Skalierung x
    \foreach \x/\label in {0/0, 2/4, 4/8, 6/12, 8/16} {
        \draw (\x,-0.1) -- (\x,0.1);
        \node[below] at (\x,-0.3) {\small \label};
    }

    % Skalierung y
    \foreach \y/\label in {0/0, 1/20, 2/40, 3/60, 4/80} {
        \draw (-0.1,\y) -- (0.1,\y);
        \node[left] at (-0.3,\y) {\small \label};
    }

    % Hilfslinien
    \draw[lightgray, very thin] (0,1) -- (8,1);
    \draw[lightgray, very thin] (0,2) -- (8,2);
    \draw[lightgray, very thin] (0,3) -- (8,3);
    \draw[lightgray, very thin] (0,4) -- (8,4);
    \draw[lightgray, very thin] (2,0) -- (2,4);
    \draw[lightgray, very thin] (4,0) -- (4,4);
    \draw[lightgray, very thin] (6,0) -- (6,4);
    \draw[lightgray, very thin] (8,0) -- (8,4);
\end{tikzpicture}
\end{center}

\noindent\textbf{b)} Berechne die Geschwindigkeit des Joggers:

\vspace{0.3em}
\hspace{1.5em} $v =$ \luecke{5cm}

\vspace{1em}

% ============================================================================
% AUFGABE 5: Bewegungen vergleichen (**) - 5 Punkte
% ============================================================================
\noindent\textbf{Aufgabe 5 \stern{**}: Bewegungen vergleichen} \punkte{5}

\vspace{0.5em}
\noindent Das Diagramm zeigt zwei Bewegungen (A und B):

\begin{center}
\begin{tikzpicture}[scale=0.65]
    % Koordinatensystem
    \draw[->] (0,0) -- (7,0) node[right] {$t$ in s};
    \draw[->] (0,0) -- (0,5) node[above] {$s$ in m};

    % Skalierung x
    \foreach \x/\label in {0/0, 1/2, 2/4, 3/6, 4/8, 5/10} {
        \draw (\x,-0.1) -- (\x,0.1);
        \node[below] at (\x,-0.2) {\small \label};
    }

    % Skalierung y
    \foreach \y/\label in {0/0, 1/10, 2/20, 3/30, 4/40} {
        \draw (-0.1,\y) -- (0.1,\y);
        \node[left] at (-0.2,\y) {\small \label};
    }

    % A: steil (v = 8 m/s → 40m in 5s)
    \draw[thick, red] (0,0) -- (2.5,4);
    \node[red, right] at (2,3.5) {\small A};

    % B: flach (v = 4 m/s → 40m in 10s)
    \draw[thick, physik] (0,0) -- (5,4);
    \node[physik, right] at (4.5,3.5) {\small B};
\end{tikzpicture}
\end{center}

\begin{enumerate}[label=\alph*), leftmargin=2em]
    \item Welche Bewegung ist schneller? \luecke{1.5cm} \hfill (1 P)
    \item Begründe deine Antwort in einem Satz:\\[0.5em]
          \luecke{13cm} \hfill (2 P)
    \item Nach welcher Zeit hat Bewegung A den Weg $s = 40$ m erreicht? $t =$ \luecke{1.5cm} \hfill (2 P)
\end{enumerate}

\vspace{1em}

% ============================================================================
% AUFGABE 6: Transfer (***) - 6 Punkte
% ============================================================================
\noindent\textbf{Aufgabe 6 \stern{***}: Bewegungsgeschichte interpretieren} \punkte{6}

\vspace{0.5em}
\noindent Das Diagramm zeigt die Bewegung von Lisa auf dem Schulweg:

\begin{center}
\begin{tikzpicture}[scale=0.55]
    % Koordinatensystem
    \draw[->] (0,0) -- (11,0) node[right] {$t$ in min};
    \draw[->] (0,0) -- (0,6) node[above] {$s$ in m};

    % Skalierung x
    \foreach \x/\label in {0/0, 2/2, 4/4, 6/6, 8/8, 10/10} {
        \draw (\x,-0.1) -- (\x,0.1);
        \node[below] at (\x,-0.2) {\small \label};
    }

    % Skalierung y
    \foreach \y/\label in {0/0, 1/100, 2/200, 3/300, 4/400, 5/500} {
        \draw (-0.1,\y) -- (0.1,\y);
        \node[left] at (-0.2,\y) {\small \label};
    }

    % Bewegung: schnell - Stillstand - langsam
    \draw[thick, physik] (0,0) -- (2,2);  % schnell (0-2 min: 0-200m)
    \draw[thick, physik] (2,2) -- (5,2);  % Stillstand (2-5 min: 200m)
    \draw[thick, physik] (5,2) -- (10,5); % langsam (5-10 min: 200-500m)
\end{tikzpicture}
\end{center}

\begin{enumerate}[label=\alph*), leftmargin=2em]
    \item Beschreibe, was in den drei Abschnitten des Diagramms passiert:

    \vspace{0.3em}
    \textbf{Abschnitt 1} (0--2 min): \luecke{10cm}

    \vspace{0.3em}
    \textbf{Abschnitt 2} (2--5 min): \luecke{10cm}

    \vspace{0.3em}
    \textbf{Abschnitt 3} (5--10 min): \luecke{10cm}
    \hfill (3 P)

    \item Erfinde eine passende Geschichte, die zu diesem Diagramm passt:\\[0.5em]
          \luecke{13cm}\\[0.8em]
          \luecke{13cm}\\[0.8em]
          \luecke{13cm} \hfill (3 P)
\end{enumerate}

\vspace{1.5em}

% ============================================================================
% PUNKTEÜBERSICHT
% ============================================================================
\hrule
\vspace{0.5em}

\begin{center}
\renewcommand{\arraystretch}{1.3}
\begin{tabular}{|l|c|c|c|}
\hline
\textbf{Niveau} & \textbf{Aufgaben} & \textbf{max. Punkte} & \textbf{erreicht} \\
\hline
\stern{*} Basis & 1, 2, 3 & 20 P & \hspace{1.2cm}/20 \\
\hline
\stern{**} Standard & + 4, 5 & 33 P & \hspace{1.2cm}/33 \\
\hline
\stern{***} Erweitert & + 6 & 39 P & \hspace{1.2cm}/39 \\
\hline
\end{tabular}
\end{center}

\vspace{0.3em}

\begin{center}
\small
\begin{tabular}{|c|c|c|c|c|c|c|}
\hline
\textbf{Note} & 1 & 2 & 3 & 4 & 5 & 6 \\
\hline
\textbf{ab Punkte} & 37 & 31 & 23 & 16 & 8 & 0 \\
\hline
\textbf{Prozent} & $\geq$95\% & $\geq$80\% & $\geq$60\% & $\geq$40\% & $\geq$20\% & $<$20\% \\
\hline
\end{tabular}
\end{center}

\end{document}
